\section{Код программ для построения множеств Мандельброта и Жюлиа}

\subsection{Множество Мандельброта}

\begin{lstlisting}
import numpy as np
import matplotlib.pyplot as plt

def mandelbrot(h, w, maxit=20):
    y, x = np.ogrid[-1.4:1.4:h * 1j, -2:0.8:w * 1j]
    c = x + y * 1j
    z = c
    divtime = maxit + np.zeros(z.shape, dtype=int)
    for i in range(maxit):
        z = z ** 2 + c
        diverge = z * np.conj(z) > 2 ** 2
        div_now = diverge & (divtime == maxit)
        divtime[div_now] = i
        z[diverge] = 2
    return divtime


plt.imshow(mandelbrot(20000, 20000))
plt.show()

\end{lstlisting}

\subsection{МножествоЖюлиа}

\begin{lstlisting}
import numpy as np
import matplotlib.pyplot as plt

def julia(h, w, maxit=100, c=0 + 0j):

    y, x = np.ogrid[-1.5:1.5:h * 1j, -1.5:1.5:w * 1j]
    z = x + y * 1j

    divtime = maxit + np.zeros(z.shape, dtype=int)

    for i in range(maxit):
        z = z ** 2 + c
        diverge = z * np.conj(z) > 2 ** 2
        div_now = diverge & (divtime == maxit)
        divtime[div_now] = i
        z[diverge] = 2

    return divtime

plt.imshow(julia(2000, 2000, maxit=100, c=-0.5251993 + 0.5251993j))
plt.show()
\end{lstlisting}

